\section*{Bab \ref{ch:g12:stereochemistry} --- Stereokimia: Kiralitas dan Isomer}

\begin{solution}{exo:g12:stereochemistry:1}
Butan-2-ol: karbon 2 (\ce{H}, \ce{OH}, \ce{CH3}, \ce{CH2CH3}).
Propan-2-ol: tidak ada (karbon 2 membawa dua \ce{CH3} yang identik).
2-Klorobutana: karbon 2. Asam laktat: karbon 2 (\ce{H}, \ce{OH},
\ce{CH3}, \ce{COOH}).
\end{solution}

\begin{solution}{exo:g12:stereochemistry:2}
Yang pertama: kedua \ce{Cl} (berprioritas atas \ce{H}) pada sisi yang
sama: \omterm{def:g12:stereochemistry:z-e}{Z}. Yang kedua: di kiri, \ce{CH3} berprioritas, sisi atas; di
kanan, \ce{CH2CH3} berprioritas, sisi bawah: \omterm{def:g12:stereochemistry:z-e}{E}.
\end{solution}

\begin{solution}{exo:g12:stereochemistry:3}
\omterm{def:g12:stereochemistry:chiral}{Kiral}: sarung tangan, sekrup. Tidak \omterm{def:g12:stereochemistry:chiral}{kiral}: sendok, kubus (masing-masing
memiliki bidang cermin). \ce{CHFClBr} \omterm{def:g12:stereochemistry:chiral}{kiral} (empat \omterm{def:g7:atoms-and-molecules:atom}{atom} berbeda pada
karbon); \ce{CH2Cl2} tidak.
\end{solution}

\begin{solution}{exo:g12:stereochemistry:4}
\omterm{def:g6:pure-substances-and-mixtures:mixture}{Campuran} dua \omterm{def:g12:stereochemistry:enantiomer}{enantiomer} dalam jumlah yang sama. Karvon rasemat
mengandung kedua \omterm{def:g12:stereochemistry:enantiomer}{enantiomer}, sehingga hidung menerima kedua bau
sekaligus.
\end{solution}

\begin{solution}{exo:g12:stereochemistry:5}
Tidak: karbon 1 membawa dua \omterm{def:g7:atoms-and-molecules:atom}{atom} yang identik, \ce{H} dan \ce{H};
menukar keduanya menghasilkan \omterm{def:g7:atoms-and-molecules:molecule}{molekul} yang sama.
\end{solution}

\begin{solution}{exo:g12:stereochemistry:6}
\setchemfig{atom sep=2em}
\chemfig{C(-[2]CH_2CH_3)(-[4]H_3C)(<[:-30]OH)(<:[:-150]H)} \quad cermin
\quad \chemfig{C(-[2]CH_2CH_3)(-[0]CH_3)(<[:-150]HO)(<:[:-30]H)}
\end{solution}

\begin{solution}{exo:g12:stereochemistry:7}
2-Klorobutana: 2 (satu \omterm{def:g12:stereochemistry:chiral}{karbon asimetris}). 2,3-Diklorobutana: 3 (dua
\omterm{def:g12:stereochemistry:chiral}{karbon asimetris} dengan belahan yang serupa: sepasang \omterm{def:g12:stereochemistry:enantiomer}{enantiomer} dan
sebuah bentuk meso). Pent-2-ena: 2 (\omterm{def:g12:stereochemistry:z-e}{Z} dan \omterm{def:g12:stereochemistry:z-e}{E}).
\end{solution}

\begin{solution}{exo:g12:stereochemistry:8}
\omterm{def:g12:stereochemistry:enantiomer}{Enantiomer}; \omterm{def:g12:stereochemistry:diastereomer}{diastereomer}; \omterm{def:g12:stereochemistry:diastereomer}{diastereomer}; \omterm{def:g7:atoms-and-molecules:molecule}{molekul} yang sama (asam
meso-tartrat dapat di\-tumpang\-tindih\-kan dengan bayangan cerminnya).
\end{solution}

\begin{solution}{exo:g12:stereochemistry:9}
Dalam \omterm{def:g12:stereochemistry:conformation}{konformasi} goyang dengan gugus \ce{CH3} yang saling
berseberangan, gugus-gugus besar itu sejauh mungkin satu sama lain:
\omterm{def:g12:stereochemistry:conformation}{konformasi} itu paling stabil. Dalam \omterm{def:g12:stereochemistry:conformation}{konformasi} eklips, kedua \ce{CH3}
saling berhadapan dan saling berdesakan.
\end{solution}

\begin{solution}{exo:g12:stereochemistry:10}
Kedua belahannya, masing-masing \ce{CH(OH)COOH}, saling bayangan
cermin: sebuah bidang di antara kedua karbon pusat membelah \omterm{def:g7:atoms-and-molecules:molecule}{molekul} itu
menjadi dua belahan yang saling bayangan cermin. \omterm{def:g7:atoms-and-molecules:molecule}{Molekul} dengan bidang
cermin adalah bayangan cerminnya sendiri: tidak \omterm{def:g12:stereochemistry:chiral}{kiral}.
\end{solution}

\begin{solution}{exo:g12:stereochemistry:11}
\setchemfig{atom sep=1.7em}
(\omterm{def:g12:stereochemistry:z-e}{Z})-pent-2-ena \chemfig{C(-[:150]H_3C)(-[:210]H)=C(-[:30]CH_2-CH_3)(-[:-30]H)}
dan (\omterm{def:g12:stereochemistry:z-e}{E})-pent-2-ena \chemfig{C(-[:150]H_3C)(-[:210]H)=C(-[:30]H)(-[:-30]CH_2-CH_3)}.
\end{solution}

\begin{solution}{exo:g12:stereochemistry:12}
Sama seperti asam tartrat, dengan \ce{Br} sebagai ganti \ce{OH} dan
\ce{CH3} sebagai ganti \ce{COOH}: kedua \ce{Br} pejal, keduanya
putus-putus (sepasang \omterm{def:g12:stereochemistry:enantiomer}{enantiomer}), serta satu pejal satu putus-putus,
yang memiliki bidang cermin: bentuk meso. Tiga \omterm{def:g12:stereochemistry:stereoisomer}{stereoisomer}.
\end{solution}

\begin{solution}{exo:g12:stereochemistry:13}
Setengahnya. Ahli kimia dapat memisahkan kedua \omterm{def:g12:stereochemistry:enantiomer}{enantiomer} (seperti
Pasteur, dengan cara lain), atau hanya membuat \omterm{def:g12:stereochemistry:enantiomer}{enantiomer} yang berguna,
dengan reaksi yang juga \omterm{def:g12:stereochemistry:chiral}{kiral} (katalis \omterm{def:g12:stereochemistry:chiral}{kiral}, enzim, bahan awal yang
\omterm{def:g12:stereochemistry:chiral}{kiral}).
\end{solution}

\begin{solution}{exo:g12:stereochemistry:14}
Tiga: (2E,4E), (2Z,4Z), dan (2E,4Z), yang terakhir merupakan \omterm{def:g7:atoms-and-molecules:molecule}{molekul}
yang sama dengan (2Z,4E) jika dibaca dari ujung yang lain.
\end{solution}

\begin{solution}{exo:g12:stereochemistry:15}
$2^3 = 8$ \omterm{def:g12:stereochemistry:stereoisomer}{stereoisomer}, yaitu 4 pasang \omterm{def:g12:stereochemistry:enantiomer}{enantiomer}. Setiap \omterm{def:g12:stereochemistry:stereoisomer}{stereoisomer}
memiliki satu \omterm{def:g12:stereochemistry:enantiomer}{enantiomer} dan $8 - 2 = 6$ \omterm{def:g12:stereochemistry:diastereomer}{diastereomer}.
\end{solution}

\begin{solution}{pb:g12:stereochemistry:1}
\textbf{1.} \ce{C4H6O6}; $M = 4 \times 12.0 + 6 \times 1.0 + 6 \times 16.0
= \qty{150.0}{g/mol}$.

\textbf{2.} Dua gugus karboksil (asam) dan dua gugus hidroksil
(\omterm{def:g11:functional-groups:alcohol}{alkohol}).

\textbf{3.} Kedua karbon pusat, masing-masing terikat pada \ce{H},
\ce{OH}, \ce{COOH}, dan \ce{CH(OH)COOH}.

\textbf{4.} $2^2 = 4$.

\textbf{5.} Keduanya dengan baji pejal.

\textbf{6.} Kedua \ce{OH} dengan baji putus-putus; bayangan itu tidak
dapat ditumpangtindihkan dengan L: itulah asam D-tartrat, \omterm{def:g12:stereochemistry:enantiomer}{enantiomernya}.

\textbf{7.} Jika diputar di sekitar ikatan tengahnya sehingga kedua
\ce{OH} menghadap ke arah yang sama, \omterm{def:g7:atoms-and-molecules:molecule}{molekul} itu memiliki sebuah bidang
di antara kedua karbon pusatnya yang mencerminkan belahan yang satu ke
belahan yang lain.

\textbf{8.} Tidak: itulah asam meso-tartrat.

\textbf{9.} Dari keempat kombinasi (pejal, pejal), (putus-putus,
putus-putus), (pejal, putus-putus), (putus-putus, pejal), dua yang
terakhir merupakan \omterm{def:g7:atoms-and-molecules:molecule}{molekul} yang sama, yaitu bentuk meso, karena kedua
belahan \omterm{def:g7:atoms-and-molecules:molecule}{molekulnya} serupa.

\textbf{10.} L dan D: \omterm{def:g12:stereochemistry:enantiomer}{enantiomer}. L dan meso: \omterm{def:g12:stereochemistry:diastereomer}{diastereomer}.

\textbf{11.} Juga pada \qty{169}{\celsius}: \omterm{def:g12:stereochemistry:enantiomer}{enantiomer} memiliki sifat
fisika yang sama.

\textbf{12.} Tidak: senyawa itu \omterm{def:g12:stereochemistry:diastereomer}{diastereomer} L, senyawa lain dengan
sifat-sifatnya sendiri.

\textbf{13.} Reseptor di hidung bersifat \omterm{def:g12:stereochemistry:chiral}{kiral} dan dapat membedakan
bayangan cermin; termometer, yang tidak \omterm{def:g12:stereochemistry:chiral}{kiral}, tidak dapat.

\textbf{14.} Tidak, \omterm{def:g6:pure-substances-and-mixtures:mixture}{campuran} itu terdiri atas dua senyawa; secara
keseluruhan \omterm{def:g6:pure-substances-and-mixtures:mixture}{campuran} itu tidak aktif optis, karena kedua \omterm{def:g12:stereochemistry:enantiomer}{enantiomernya}
saling meniadakan.

\textbf{15.} Hanya satu \omterm{def:g12:stereochemistry:enantiomer}{enantiomer}: setiap kristal hanya tersusun atas L
atau hanya atas D.

\textbf{16.} Masing-masing \qty{5.0}{g}.

\textbf{17.} Bentuk meso adalah senyawa lain, bukan bagian dari \omterm{def:g12:stereochemistry:enantiomer}{campuran
rasemat} L dan D.

\textbf{18.} Dengan melarutkan setiap tumpukan: larutan-larutannya
memutar cahaya terpolarisasi dengan besar yang sama ke arah yang
berlawanan (dan kristal-kristalnya saling bayangan cermin).

\textbf{19.} Empat: kedua \omterm{def:g12:stereochemistry:chiral}{karbon asimetrisnya} membawa gugus yang berbeda
(\ce{Cl} pada yang satu, \ce{OH} pada yang lain), sehingga tidak ada
kombinasi yang memiliki bidang cermin, dan tidak ada dua kombinasi yang
merupakan \omterm{def:g7:atoms-and-molecules:molecule}{molekul} yang sama.

\textbf{20.} Tiga: L, D, dan meso.
\end{solution}
